% در این فایل، عنوان پایان‌نامه، مشخصات خود، متن تقدیمی‌، ستایش، سپاس‌گزاری و چکیده پایان‌نامه را به فارسی، وارد کنید.
% توجه داشته باشید که جدول حاوی مشخصات پایان‌نامه/رساله و همچنین، مشخصات داخل آن، به طور خودکار، درج می‌شود.
%%%%%%%%%%%%%%%%%%%%%%%%%%%%%%%%%%%%

\thispagestyle{empty}

% \includepdf[pages=-]{firstpages.pdf}


% دانشگاه خود را وارد کنید
\university{تربیت مدرس}
% دانشکده، آموزشکده و یا پژوهشکده  خود را وارد کنید
\faculty{علوم ریاضی}
% گروه آموزشی خود را وارد کنید
\degree {کارشناسی ارشد علوم کامپیوتر گرایش داده کاوی}
% گروه آموزشی خود را وارد کنید
\subject{علوم کامپیوتر }
% گرایش خود را وارد کنید
\field{علوم کامپیوتر (داده کاوی)}
% عنوان پایان‌نامه را وارد کنید
\title{بهبود مدل‌های انتشار مبتنی بر مبدل‌های بینایی برای تولید تصویر}
% نام استاد(ان) راهنما را وارد کنید
\firstsupervisor{دکتر منصور رزقی آهق}
%\secondsupervisor{استاد راهنمای دوم}
% نام استاد(دان) مشاور را وارد کنید. چنانچه استاد مشاور ندارید، دستور پایین را غیرفعال کنید.
%\firstadvisor{نام استاد(دان) مشاور}
%نام استاد مشاور را وارد کنید. چنانچه دو استاد مشاور دارید  دستور پایین را نیز فعال کنید.
%\secondadvisor{نام استاد(دان) مشاور}
% نام پژوهشگر را وارد کنید
\name{مصطفی}
% نام خانوادگی پژوهشگر را وارد کنید
\surname{شهبازی دیل}
% تاریخ پایان‌نامه را وارد کنید
\thesisdate{بهمن ۱۴۰۴}
% کلمات کلیدی پایان‌نامه را وارد کنید
%\keywords{داده‌های بقا فضایی، مدل‌های شکنندگی فضایی}
% کلمات کلیدی پایان‌نامه را وارد کنید
\keywords{مبدل‌های بینایی، مدل‌های انتشار، ویژگی‌های محلی و سراسری، تجمیع ویژگی، تنظیم دقیق پارامتر-کارا}
\newpage
%\fa-abstract{\noindent
%در

%}
%\newpage
\thispagestyle{empty}
\vtitle
%\newpage
%\thispagestyle{empty}
%\clearpage
%~~~
%\newpage
%\thispagestyle{empty}
%\input{rights}
%\newpage
%\thispagestyle{empty}
%\clearpage
%~~~
%\newpage
%\thispagestyle{empty}
%\centerline{{\includegraphics[width=20 cm]{replyrecord}}}
%\newpage
%\thispagestyle{empty}
%\clearpage
%~~~
% \newpage
%  % پایان‌نامه خود را تقدیم کنید!
% \begin{acknowledgementpage}

% \vspace{4cm}

% {\nastaliq
% {\Large
%  تقدیم به
% \vspace{1.5cm}

% \newdimen\xa
% \xa=\textwidth
% \advance \xa by -14cm
% \hspace{\xa}
% پدر و مادر مهربانم
% \vspace{1.5cm}

% \newdimen\xa
% \xa=\textwidth
% \advance \xa by -12cm
% \hspace{\xa}
% %و زنان آزاده کشورم
% \vspace{1.5cm}

% \newdimen\xa
% \xa=\textwidth
% \advance \xa by -10cm
% \hspace{\xa}

% \vspace{1.5cm}

% \newdimen\xa
% \xa=\textwidth
% \advance \xa by -8cm
% \hspace{\xa}
% }}
% \end{acknowledgementpage}
%\newpage
%\thispagestyle{empty}
%\clearpage
%~~~
%%%%%%%%%%%%%%%%،%%%%%%%%%%%%%%%%%%%%
%\newpage
%\thispagestyle{empty}
% ستایش
%\baselineskip=.750cm
%\ \\ \\
%
%{\nastaliq
%رسيدن، به دانش است و به كردار نیک...%
%}\\
%\vspace{.5cm}\\
%{\scriptsize\nastaliq
%{
% و بی دانش به كردار نيك هم نتوان رسيد، كه نيكى را پيشتر ببايد شناختن، آنگاه بجاى آوردن. پس دانش به همه حال مى‌ببايد تا به رستگارى %توان رسيدن. و چون دانش راه آمد، به بهترين چيزها كه آدمى را تواند بودن. و در اوّل آفرينش حاصل نيست و بعضى از آن بى‌رنج و انديشه حاصل شود، پس هرآينه مهمتر چيزى باشد كه در حاصل كردنش عمر گذرانند، ليكن برخى هست كه بى‌انديشه حاصل آيد و بعضى را ناچار به انديشه حاجت بود، و آنچه به انديشه حاصل شود دانسته‌اى خواهد كه درو انديشه كنند تا اين نادانسته بدان انديشه كه در آن دانسته كنند دانسته شود، و از هر دانسته هر نادانسته را نتوان شناخت، بلكه هر نادانسته را به دانسته‌اى كه در خور او بود توان شناخت. و منطق آن علم است كه درو راه انداختن نادانسته به دانسته دانسته شود...
% }}
%
%\vspace{.5cm}
%{\nastaliq
%\newdimen\xb
%\xb=\textwidth
%\advance \xb by -8.5cm
%\hspace{\xb}
%پس منطق ناگزير آمد بر جوينده‌ی رستگاری.
%\RTLfootnote{مقدمه‌ی رساله‌ی منطق دانشنامه‌ی علائی، شیخ‌الرئیس ابن‌سینا}
%}
%\newpage
%\thispagestyle{empty}
%\clearpage
%~~~
%%%%%%%%%%%%%%%%%%%%%%%%%%%%%%%%%%%%
% \newpage
% \thispagestyle{empty}
% %قدردانی
% {\nastaliq
% قدردانی
% }
% \\[2cm]
% {\scriptsize\nastaliq
% منّت خدای را عزّ و جل که طاعتش موجب قربت است و به شکر اندرش مزید نعمت. هر نفسی که فرو می‌رود مُمِدّ حیات است و چون بر می‌آید مُفَرّح ذات. پس در هر نفسی دو نعمت موجود است و بر هر نعمتی شکری واجب(دیباچه گلستان).
% }


% % با استفاده از دستور زیر، امضای شما، به طور خودکار، درج می‌شود
% \signature
% %
\newpage
\thispagestyle{empty}
\baselineskip=1.0cm
\ \\ \\
\begin{large}
\begin{center}
چکیده
\end{center}
\end{large}
%}}
%\\[4cm]
مدل‌های انتشار مبتنی بر مبدل، مانند \lr{DiT}، کیفیت بالایی در تولید تصویر شرطی ارائه می‌کنند، اما تنظیم دقیق کامل آن‌ها برای دامنه‌ها یا نیازهای جدید از نظر حافظه، زمان آموزش و ذخیره‌سازی پرهزینه است. این پژوهش معماری \lr{DuoDiT} را برای تنظیم دقیق پارامتر-کارای مدل‌های \lr{DiT} معرفی می‌کند. در این معماری، بدنه اصلی \lr{DiT-XL/2} منجمد باقی می‌ماند و یک شاخه دوم با اندازه پچ کوچک‌تر، ویژگی‌های محلی و ریزدانه را استخراج می‌کند. برای هم‌تراز کردن این شاخه با توکن‌های بدنه اصلی، از مکانیزم توکن کلاس محلی استفاده می‌شود که هر گروه مکانی از توکن‌های ریزدانه را به‌صورت محتوا-آگاه خلاصه می‌کند. ارزیابی اصلی روی ۵۰ هزار عکس تولیده به صورت شرطی روی مجموعه داده 
\lr{ImageNet-1K}
 در وضوح 
 $256 \times 256$
   انجام شده است. نسخه 
   \lr{DuoDiT-200}
    با 
    \lr{17.95M}
     پارامتر آموزش‌پذیر از مجموع \lr{690.39M} پارامتر، به \lr{FID=2.219}، \lr{KID=0.0014} و \lr{LPIPS-Mean=0.7162} می‌رسد. همچنین در مطالعه انسانی، نمونه‌های \lr{DuoDiT} در \lr{53.26\%} مقایسه‌ها نسبت به \lr{LightningDiT} ترجیح داده شدند. نتایج نشان می‌دهد که افزودن یک مسیر سبک و وضوح‌بالا می‌تواند بدون آموزش مجدد کامل مدل، کیفیت ادراکی و کارایی تنظیم دقیق را بهبود دهد.\\
\textit{ \textbf{ واژه‌های کلیدی:}}
مبدل‌های بینایی، مدل‌های انتشار، ویژگی‌های محلی و سراسری، تجمیع ویژگی، تنظیم دقیق پارامتر-کارا.
%\newpage
%\thispagestyle{empty}
%\clearpage
%~~~
%\newpage
%{\small
%\abstractview
%}
%\newpage
%\thispagestyle{empty}
%\clearpage
%~~~
\newpage
